Visualising the code review process has been an incredibly educational process and that developed skills in several areas:

\begin{enumerate}
    \item Problem-solving —
    there were several occasions where an error occurred in the code of the project and a solution had
    to be formulated with limited time and resources.
    The ability to look at a problem
    and systematically work out what was going wrong became a regular occasion in this project,
    some of which took multiple days at a time.
    Patience and vigilance were a virtue for this purpose.
    \item Organisation — throughout the entire project there was a strict timeline in place;
    milestones were allocated usually a week of time each to complete
    which had to be adapted to for the requirements to be met on time.
    \item Communication — the Scrum development framework was followed in this project,
    and offered valuable experience in daily stand-up meetings
    that allowed for voicing of any concerns or issues faced throughout the project.
    \item Support — the project was never going to be straight forward,
    and additional support was inevitable due to the complex nature of what was being made.
    The ease
    in which support could be obtained could not have been made any easier and simpler with the Scrum methodology alongside the constant back-and-forth between members of the team.
    \item Knowledge — the wealth of knowledge that was obtained from completing a project of this calibre was staggering.
    This could range from programming skills in JavaScript and TypeScript,
    all the way to how a full-stack application is built from the ground-up.
    These abilities will become cornerstones of future career opportunities,
    since they are so versatile and applicable in many areas regardless of the specific project.
    \item Connections — outside software engineering itself,
    the project also gave way for many connections to be formed, and not just academically.
    A paper will also be written on the project,
    an accomplishment
    that will end up being repeated in future ventures,
    so being able to do it competently adds a layer of insight that,
    otherwise, would not have been achievable.
\end{enumerate}

In conclusion,
I would like to thank DongGyun Han
for giving me the opportunity to research his work and trusting me to complete it to his standards.
His advice and teachings will be forever valuable in my future projects and work that I complete.
I am very proud of the final product that we have created and, without his support, nothing would have been possible.
Furthermore,
I would also like to thank Matteo Sammartino for accepting me into the UROP program despite my lack of experience.
The decision to take on this project is definitely not one I will ever regret.